\section{A formal detector}
\label{sec:detector}

We first separate the candidate from the relation space.  Throughout this
section the genus \(g\) is fixed as in Definition~\ref{def:candidate},
but the number of labels is allowed to vary, including zero.
No Chow relation on a positive-genus vertex will be imposed in the
construction.

\begin{proposition}
\label{prop:detector}
There are a supercommutative \(\Q\)-algebra \(R\), a finite free
supermodule \(V\) over \(R\), and \(\Q\)-linear maps
\[
 t_m:\Salg_{g,m}\longrightarrow V^{\otimes_R m}
 \qquad(m\geq 0)
\]
such that \(t_m(\Pspan_{g,m})=0\) for every \(m\), whereas
\(t_n(Y)\neq0\).  In particular \(Y\notin\Pspan_{g,n}\).
\end{proposition}

The maps are tests on formal strata.  In particular, they are not
asserted to factor through the Chow ring.

\subsection{An isotropic cubic Frobenius algebra}

The state space has the curve pairing familiar from Frobenius-algebra
and CohFT constructions \cite{Abrams1996,PZ2019}. The exterior
coefficient ring below allows an independent cubic on every triple of
colors. We check associativity and the handle formula directly, since
these particular identities govern the graph test.

Let
\[
 R=\Lambda_{\Q}\bigl(\theta_{ijk}:1\leq i<j<k\leq g\bigr),
\]
with each generator odd.  Let \(W\) be a purely odd symplectic vector
space of dimension \(2g\), with complementary Lagrangians
\(W_+\) and \(W_-\).  Choose bases \(a_1,\ldots,a_g\) and
\(b_1,\ldots,b_g\) such that
\(\langle a_i,b_j\rangle=\delta_{ij}\).
Extend the symbols \(\theta_{ijk}\) alternately to all triples of
indices.  There is an alternating, odd-coefficient three-form
\(\Phi\) on \(W\) defined by
\[
 \Phi(a_i,a_j,a_k)=\theta_{ijk},
 \qquad
 \Phi(w_1,w_2,w_3)=0
 \quad\text{if some }w_i\in W_-.
\]

Consider the free supermodule
\[
 V=R\otimes_{\Q}\bigl(\Q1\oplus W\oplus\Q e\bigr),
\]
where \(1,e\) are even.  Its trace, denoted by \(\int\), takes \(e\)
to \(1\) and is zero on \(1,W\).  Give \(V\) the even pairing which
is hyperbolic on \(\Q1\oplus\Q e\), the prescribed symplectic pairing
on \(W\), and zero between these two summands.  Declare \(1\) to be
the unit, \(e^2=eW=0\), and determine the remaining products by
\[
 \int w_1w_2=\langle w_1,w_2\rangle,
 \qquad
 \int w_1w_2w_3=\Phi(w_1,w_2,w_3).
\]
Equivalently, the product of two standard odd vectors is its
symplectic pairing times \(e\), plus the vector dual to the cubic
functional in the third input.  The latter vector has odd
coefficients and belongs to \(R\otimes W_-\).

Write \(\mu\) for multiplication and \(\Delta\) for its adjoint
coproduct.

\begin{lemma}
\label{lem:detector-algebra}
These rules give a supercommutative associative Frobenius algebra.
The product of three standard odd vectors belongs to \(Re\), and
the product of four such vectors is zero.  Its handle element is
\[
 \mu\Delta(1)=(2-2g)e.
\]
\end{lemma}

\begin{proof}
The pairing is nondegenerate, so the rules determine the product.
Supercommutativity follows from alternation of the symplectic form
and of \(\Phi\), with coefficients moved using the super sign rule.
For example,
\[
 a_i a_j=-\sum_k\theta_{ijk}b_k,
 \qquad a_i b_j=\delta_{ij}e,
\]
and therefore \((a_i a_j)a_k=\theta_{ijk}e
=a_i(a_j a_k)\).

More generally, pair the associator of three vectors from \(W\)
with a fourth basis vector.  Pairing with \(1\) gives the cyclic
symmetry of the alternating cubic, and pairing with \(e\) gives
zero.  If the fourth vector is in \(W\), both expressions are
contractions of two pair-product tensors.  Their \(W\)-channel
contractions vanish because the cubic is supported on a
Lagrangian; their hyperbolic-channel contractions vanish because
\(e^2=eW=0\).  Nondegeneracy proves associativity.  Cases involving
\(1\) or \(e\) are immediate.  The same product description proves
the assertions about three and four odd inputs.

The hyperbolic channels contribute \(2e\) to the handle element,
and the \(2g\) odd channels contribute \(-2ge\).  The cubic
perturbation has zero internal trace, because a symplectic
contraction pairs the two different Lagrangians.  This gives the
displayed formula.
\end{proof}

Set
\[
 \lambda=(2g-2)e.
\]
For a nonempty finite label set \(L\), write
\(\Delta_L:V\to V^{\otimes L}\) for the coproduct adjoint to
iterated multiplication; for a singleton it is the identity.
Coproducts are coassociative and satisfy the projection formula.
In particular, writing \(\Delta_{ij}(1)\) with units in all other
slots, we have
\begin{align}
 \Delta_{ij}(1)\Delta_{ik}(1)&=\Delta_{ijk}(1), \label{eq:detector-diagonals}\\
 \Delta_{ij}(1)^2&=-\Delta_{ij}(\lambda). \label{eq:detector-self}\\
 \int_j\Delta_{ij}(v)&=v. \label{eq:detector-counit}
\end{align}
The minus sign in \eqref{eq:detector-self} will match the normal
bundle of a rational-tail boundary divisor.

\subsection{Compression of rational tails}

Call a graph a \emph{rational-tails graph} if it is a tree with one
genus-\(g\) vertex, called its root, and all other vertices rational.
Every branch away from the root has nonempty labeled descendants.
Such a graph has no automorphisms fixing the labels: the descendant
label sets distinguish its branches, recursively.

We initially work without the degree cutoff, discard all graphs
which are not rational-tails graphs, and allow genuine genus-zero
Chow identities on rational parameter factors.  This is only an
auxiliary computation; we will prove that its output respects the
cutoff in \(\Salg_{g,m}\). We use the genus-zero Chow boundary
presentation of \cite[Section~4, Theorem~1]{Keel1992}.

For each label \(i\), put
\[
 \bar\psi_i=\psi_i-\sum_{\substack{T\ni i\\ |T|\geq2}}D_T.
\]
If the labels are ordered \(1,\ldots,m\), let \(L_x\) be the
pullback of \(\psi_x\) from the space retaining just
\(1,\ldots,x\).  Thus
\[
 L_x=\psi_x-
 \sum_{\substack{T\ni x,\ |T|\geq2\\T\setminus\{x\}\subseteq\{x+1,\ldots,m\}}}
 D_T.
\]
For \(j\geq1\), define
\[
 \bar\kappa_j=\kappa_j-\sum_{x=1}^m L_x^j.
\]
The leading, graphless change of variables is
\(\bar\psi_i=\psi_i\) and
\(\bar\kappa_j=\kappa_j-\sum_i\psi_i^j\), which is an invertible
polynomial change.  Hence the bar monomials form a complement to
the boundary subspace before truncation.

Define \(t_m\) on that complement by
\begin{equation}
 t_m\left(\prod_i\bar\psi_i^{d_i}
                 \prod_{j\geq1}\bar\kappa_j^{c_j}\right)
 =
 \begin{cases}
   \lambda_1^{d_1}\cdots\lambda_m^{d_m},
      &\text{if every }c_j=0,\\
   0,&\text{otherwise}.
 \end{cases}
 \label{eq:detector-prescription}
\end{equation}
Here the tensor on the right has a unit in every unused slot.
For a nontrivial rational-tails graph, integrate each non-root
decoration on its genus-zero parameter space, retaining only its
top-degree integral.  Apply \(t_a\) to the decorated root with its
\(a\) incident slots.  Finally duplicate each root slot into the
descendant labels of its branch by \(\Delta_L\).
Every branch attached to the root has at least two descendant
labels, so \(a<m\); this is a recursive definition.  For \(m=0\),
the empty tensor is \(R\), the unit maps to \(1\), and positive
bar-kappa monomials map to zero.

\begin{lemma}
\label{lem:detector-cutoff}
The prescription is well defined modulo genuine relations on
rational parameter factors and induces a map on \(\Salg_{g,m}\).
For a homogeneous expression of degree \(r\), its image is a sum
of tensors indexed by partitions of the labels into nonempty
blocks, of the form
\begin{equation}
 \bigotimes_{L}\Delta_L(\lambda^{d_L}),
 \qquad
 r=\sum_L\bigl(|L|-1+d_L\bigr).
 \label{eq:detector-blocks}
\end{equation}
In particular \(t_m\) is zero in degrees \(r>m\).
\end{lemma}

\begin{proof}
Grafting a genus-zero graph amounts to composing coproducts and
multiplying the corresponding integrals.  Coassociativity shows
that the result is independent of the internal decomposition of
a rational parameter class.  In particular a zero class on such
a factor contributes zero.

For a rational subtree with \(k\) descendant labels, top
decoration degrees together with its edges, including the edge
attaching it to the root, add to \(k-1\).  Formula \eqref{eq:detector-blocks} follows by
induction, starting from the bar prescription.  Since
\(\lambda^2=0\), a nonzero term has \(d_L\leq1\), and hence \(r\leq m\).

An excessive decoration on a rational vertex has zero integral.
At a root with \(a\) incident slots, a decoration exceeding
\(3g-3+a\) also exceeds \(a\), and its recursive test is zero.
This applies since \(g\geq2\).  Finally, under refinement the sum
of the dimensions of the pieces of a vertex decreases by the
number of new edges.  Splitting a decoration additively between
those pieces cannot make an excessive total degree admissible
on every piece.  Thus discarded terms also remain invisible
under multiplication and refinement.  Forgetting operations
require a separate argument, supplied below.
\end{proof}

\subsection{Boundary operations and the test}

We first show how the test treats multiplication by bar classes and
rational-tail divisors.  We then prove the forgetting identities
that will reduce the number of labels in the relation argument.
All tensor products, contractions, and permutations below use
the super sign convention.

\begin{lemma}[Boundary operations]
\label{lem:detector-operations}
The maps \(t_m\) are equivariant for relabeling.  For every
strata expression \(b\), they satisfy
\begin{align}
 t(\bar\psi_i b)&=\lambda_i t(b),
 &t(\bar\kappa_j b)&=0, \label{eq:detector-bars}\\
 t(\bar D_{ij}b)&=\Delta_{ij}(1)t(b). \label{eq:detector-pair}
\end{align}
More generally, gluing along \(D_T\) integrates the entire
genus-zero input and duplicates the root-side attaching slot
into the labels of \(T\).
\end{lemma}

\begin{proof}
We first establish the boundary restrictions and the graph
compatibilities used in the calculation. The descendant and kappa
identities are the forgetting and gluing formulas in
\cite[Equations~(1.7)--(1.10)]{ArbarelloCornalba1996}.  At the gluing for
\(D_T\), let \(p\) be the root-side flag and \(p'\) the
rational-side flag.  The normal class is
\(-\psi_p-\psi_{p'}\).  On
\(\Mbar_{0,T\cup\{p'\}}\), the sum of boundary splits separating
\(i\) from \(p'\) equals \(\psi_i+\psi_{p'}\).  The sum whose
side away from \(p'\) contains two fixed labels \(i,j\) equals
\(\psi_{p'}\).  These are the genus-zero cotangent boundary
formulas; they also follow by successively forgetting to three
labels and applying the cotangent pullback formula.

Consequently, \(\bar\psi_i\) restricts to the modified cotangent
of its ancestor slot on the root side, using \(p\) if \(i\in T\).
The restriction of \(\bar D_{ij}\) is the analogous pair class
on the root side, except when both labels lie in \(T\), when it
is \(-\bar\psi_p\).

The class \(\bar\kappa_j\) is independent of the chosen ordering.
Indeed, for two consecutive labels made last, write
\(D=D_{\{i,j\}}\).  The identities
\(\psi_i|_D=\psi_j|_D=0\) and
\(D|_D=-\psi_p\) give, for \(d\geq1\),
\begin{equation}
  (\psi_i-D)^d-\psi_i^d=-D\psi_p^{d-1}
  =(\psi_j-D)^d-\psi_j^d.
 \label{eq:detector-order}
\end{equation}
After pulling to the remaining labels, this is precisely the
invariance under a consecutive transposition.  To restrict
\(\bar\kappa_j\) to \(D_T\), put the labels of \(T\) last.
The kappa class splits additively between the two vertices.
The first label of \(T\) gives the root-side cotangent at \(p\);
the remaining partial cotangents on the rational factor cancel
its kappa class by the forgetful kappa formula.  The result is
exactly \(\bar\kappa_j\) on the root side.

These restrictions and the coproduct projection formula prove
\eqref{eq:detector-bars} recursively, beginning with \eqref{eq:detector-prescription}.
For \eqref{eq:detector-pair}, only the two-label tail has a
nonzero undecorated integral.  If \(i,j\) are in different
branches, the root-side pair class joins their ancestor slots.
If they are in the same branch, its restriction is
\(-\bar\psi_p\), and the required identity is
\eqref{eq:detector-self}.  This proves the assertion for an
arbitrary boundary expression as well as for a bar monomial.
The gluing assertion follows directly by coassociativity and
the factorization of genus-zero integrals.

These calculations take place in the auxiliary formal graph space
modulo relations on rational parameter factors.  Simultaneous
compatible splits produce the same tree independently of their
order, with one normal factor for every repeated edge.  Hence the
polynomial multiplications used here commute and associate.
A genuine zero on a rational factor remains zero: restriction
there is a genuine genus-zero operation, and operations at another
factor retain that zero with relabeled flags.
Lemma~\ref{lem:detector-cutoff} permits the dimension truncations.
Graphs outside the rational-tails locus remain outside it under
refinement, so their products contribute zero as well.

Finally, independence of the bar-kappa ordering, the restriction
formulas, and the symmetric coproduct construction give
equivariance under arbitrary relabeling.
\end{proof}

\subsection{Forgetting markings}

The forgetting identities are needed in both directions: pullback
will remove a zero-index leg from the Pixton recipe, and pushforward
will test a tensor slot against the two even basis vectors.

\begin{lemma}[Forgetting operations]
\label{lem:detector-forgetting}
For every strata expression \(b\), the maps \(t_m\) satisfy
\begin{equation}
 t(\pi_y^*b)=t(b)\otimes1,
 \qquad
 t(\pi_{y*}b)=\int_y t(b).
 \label{eq:detector-forget}
\end{equation}
Here \(\pi_y\) forgets the label \(y\); the identities hold
in any order of the remaining labels.
\end{lemma}

\begin{proof}
We first calculate before the degree cutoff, modulo genuine
genus-zero relations, and check at the end that every discarded
term has zero test after forgetting.  No relation on a
positive-genus parameter factor enters the calculation.

\paragraph{Formal pullback and projection.}
Forget pullback adds
\(y\) at every possible vertex.  At that vertex it replaces
\begin{equation}
 \psi_h\longmapsto\psi_h-D_{\{h,y\}},
 \qquad
 \kappa_j\longmapsto\kappa_j-\psi_y^j.
 \label{eq:detector-vertex-pull}
\end{equation}
A tail divisor pulls back as \(D_T+D_{T\cup\{y\}}\).
Against a tree, an incompatible split remains incompatible;
a compatible new edge splits a vertex.  If the split edge
already occurs, the two pullback terms give its pulled normal
class: adding \(y\) at an endpoint adds the boundary term
cutting off that half-edge together with \(y\).  Thus pullback
commutes with multiplication by the indicated polynomial
classes.

The push-projection calculation is equally local.  At a stable
vertex, put
\[
 X_i=\pi_y^*\psi_i,\qquad k_j=\pi_y^*\kappa_j,\qquad
 d_i=D_{\{i,y\}},
\]
so that \(\psi_i=X_i+d_i\) and
\(\kappa_j=k_j+\psi_y^j\).  The relevant rules are
\begin{equation}
 d_i d_l=0\ (i\neq l),\qquad
 d_i^2=-X_i d_i,\qquad
 \psi_y d_i=0,
 \label{eq:detector-forget-products}
\end{equation}
and
\begin{equation}
 \pi_{y*}1=0,\qquad
 \pi_{y*}\psi_y^d=\kappa_{d-1}\ (d\geq1),\qquad
 \pi_{y*}d_i=1.
 \label{eq:detector-forget-push}
\end{equation}
They prove projection with graphless cotangent and kappa
polynomials by substitution.  For a tail multiplier, its two
pullback splits distribute the \(\kappa_{d-1}\) in \eqref{eq:detector-forget-push}
additively over the two vertices; for \(d_i\), exactly the
compatible split pushes to the original tail.  On a tree with
a stable \(y\)-vertex this reduces to the same vertex
calculation.  If forgetting contracts a three-point vertex
with a surviving terminal label, the compatible split places
\(y\) with that label.  If it contracts a bridge, the two
normal contributions add to the normal of the combined edge.
All other compatible splits lift uniquely.  This proves the
needed push projection on boundary graphs as well.

Forgetting also preserves the rational-factor relations used in
the auxiliary computation: its pullbacks and pushforwards on
those factors are actual genus-zero operations.  Operations at
another factor retain a rational-factor zero with relabeled flags.

\paragraph{Pullback of the test.}
Both bar classes
pull to themselves.  On a boundary graph, adding \(y\) to a
rational factor increases its dimension without increasing
the pulled decoration degree, so its integral is zero.
Adding \(y\) at the root uses the smaller recursive pull
identity.  This proves the pull assertion.

\paragraph{Pushforward of the test.}
For pushforward on a boundary graph, forgetting on a rational
factor commutes with its integral; a three-point contraction
uses the counit identity \eqref{eq:detector-counit}.
Forgetting at the root again uses the smaller recursive
assertion.  On the bar-monomial complement, projection reduces
the question to powers of \(\bar\psi_y\), since all other bar
classes are pulled back.  The zeroth power pushes to zero.
For \(d>0\), telescope
\begin{equation}
 \psi_y^d-\bar\psi_y^d
 =
 \sum_{\varnothing\neq S}
 D_{S\cup\{y\}}\sum_{a=0}^{d-1}
        \psi_y^{d-1-a}\bar\psi_y^a.
 \label{eq:detector-telescope}
\end{equation}
On the rational tail with labels \(S\cup\{y\}\), the dimension
is \(|S|-1\), and only the term \(a=d-|S|\) integrates.
Its integral is
\(\int_{\Mbar_{0,|S|+2}}\psi_y^{|S|-1}=1\).
After pushforward and application of the test, \eqref{eq:detector-telescope} is
therefore
\begin{equation}
 \sum_{1\leq |S|\leq d}\Delta_S(\lambda^{d-|S|}),
 \label{eq:detector-tail-sum}
\end{equation}
with units in all unused slots.

On the other hand,
\begin{equation}
 t_m(\kappa_{d-1})
 =
 \int\lambda^d+
 \sum_{1\leq |S|\leq d}\Delta_S(\lambda^{d-|S|}).
 \label{eq:detector-kappa}
\end{equation}
For \(d=1\) this is the scalar identity
\(\kappa_0=2g-2+m\).
For \(d>1\), expand
\(\kappa_{d-1}=\bar\kappa_{d-1}+\sum_xL_x^{d-1}\).
Before pulling from the stage with labels \(1,\ldots,x\),
write the exact polynomial identity
\begin{equation}
 \psi_x^{d-1}-\bar\psi_x^{d-1}
 =
 \sum_{\substack{T\subseteq\{1,\ldots,x\}\\x\in T,\ |T|\geq2}}
 D_T\sum_{a=0}^{d-2}
       \psi_x^{d-2-a}\bar\psi_x^a .
 \label{eq:detector-partial-telescope}
\end{equation}
On the tail for \(T\), the dimension is \(|T|-2\).
Thus its integral selects exactly \(a=d-|T|\), when that
exponent lies between \(0\) and \(d-2\), and the integral
of \(\psi_x^{|T|-2}\) is \(1\).  The remaining bar cotangent
is the root-side cotangent, so this term contributes
\(\Delta_T(\lambda^{d-|T|})\), with coefficient \(+1\).
The term \(\bar\psi_x^{d-1}\) gives the singleton \(T=\{x\}\).
Pulling through the later labels adds only unit tensor
slots, by the pull identity already proved.  Every block
therefore occurs once, indexed by its largest label, with
no binomial coefficient or normal self-intersection sign.
Subtracting \eqref{eq:detector-tail-sum}
from \eqref{eq:detector-kappa} proves
\(t_m(\pi_{y*}\bar\psi_y^d)=\int\lambda^d\), as required.

\paragraph{Compatibility with the cutoff.}
A non-rational-tails graph stays outside the test under forgetting:
its nontrivial unmarked stabilization persists.

We must also check forgetting on terms discarded by the
vertexwise cutoff; this does not follow from the refinement
dimension argument alone.  An excessive decoration on a
rational factor is an actual zero in its genus-zero Chow
ring.  Its genuine forget pullback or pushforward is still
zero, including the integral used when a three-point factor
contracts.  Operations on another factor preserve this
rational-factor zero.
For an excessive root decoration of degree \(d\) with
\(a\) root slots, we have \(d\geq3g-2+a\).  Adding a mark
at that root gives an expression of degree \(d>a+1\),
so the block bound kills its test with \(a+1\) slots.
Forgetting a root mark gives degree \(d-1>a-1\), killed
by the same bound with \(a-1\) slots.  The root stays stable
because \(g\geq2\).  Adding or forgetting on a rational
factor leaves the root decoration unchanged; a contracted
rational three-point vertex merely replaces an attaching
flag without changing the root arity.
Thus discarded terms have zero test before and after each
forget operation.  Together with
Lemma~\ref{lem:detector-cutoff}, this permits every cutoff
in the displayed test identities.  It does not assert that
untruncated pullback descends unchanged to the truncated
formal algebra.
Equivariance from Lemma~\ref{lem:detector-operations} allows
the forgotten label to be placed last in any chosen order.
\end{proof}

\subsection{The zero-leg identity}

We require one precise identity for the recipe in
Definition~\ref{def:pixton-span}.  Package the partition
insertions with formal variables \(s_\ell\), one for each
allowed positive part, using exponential coefficients for
repeated parts, and put
\[
 H(T,\zeta)=\widehat C_0(T,\zeta)
             -\sum_\ell s_\ell\widehat C_\ell(T,\zeta).
\]
The variables \(s_\ell\) are purely formal; all subsequent
coefficient extractions involve finitely many of them. The logarithm
below is expanded jointly in \(T\) and the \(s_\ell\), so its constant
term at \(T=s_\ell=0\) is defined even though \(H(0,\zeta)\) depends
on the partition variables.
The cycle exponential formula for \(\widehat\kappa_\Gamma\)
rewrites the vertex characteristic as
\begin{equation}
 \exp\{-\log H\}_{\kappa,\zeta_v}.
 \label{eq:detector-characteristic}
\end{equation}
Here a monomial \(T^d\zeta^\epsilon\) in \(-\log H\) is
replaced by \(\kappa_dT^d\zeta_v^\epsilon\), including \(d=0\).
For the edges use \(U(x,x')=T\Delta_e\), so the coefficient
of total degree \(T^r\) is the degree-\(r\) recipe.
Denote this packaged series, with its usual parity
extractions, by \(\mathcal R(s;a)\).

\begin{lemma}
\label{lem:detector-unit}
If all preceding leg indices \(a_j\) are zero or one, then,
in the untruncated rational-tails computation,
\begin{equation}
 \mathcal R(s;a,0)
 -\sum_\ell s_\ell\mathcal R(s;a,\ell)
 =\pi_y^*\mathcal R(s;a).
 \label{eq:detector-unit}
\end{equation}
\end{lemma}

\begin{proof}
The kappa pullback formula, including
\(\kappa_0^{\mathrm{old}}=\kappa_0^{\mathrm{new}}-1\),
multiplies the new characteristic \eqref{eq:detector-characteristic} by
\(H(\psi_yT,\zeta_v)\).  This is exactly the linear
combination of new leg factors on the left.

It remains to check bubbles introduced in pulling the old
leg and half-edge cotangents.  For a function \(f(x)\), this
pull gives the boundary correction
\((f(0)-f(x))/x\), where \(x\) is the stable-side cotangent.
At a new rational three-point vertex, its characteristic is
\(H(0,\epsilon)^{-1}\); the new \(H\)-leg cancels it.
Let \(\eta,\eta'\) be stable-side parity variables, and write
\[
 A_x=A(\eta xT),\quad B_x=B(\eta xT),
 \qquad
 A_{x'}=A(\eta'x'T),\quad B_{x'}=B(\eta'x'T).
 \]
The edge adjoining that bubble has factor
\begin{equation}
 \frac{\eta(B_x+1)+\epsilon(1-A_x)}{x}.
 \label{eq:detector-bubble-edge}
\end{equation}
If the bubble carries an old leg, its constant factor is
\(1\) or \(-\epsilon\).  Extracting \(\epsilon\) in \eqref{eq:detector-bubble-edge}
then gives \((1-A_x)/x\) or
\(-\eta(B_x+1)/x\), respectively.  These are exactly the
two divided differences.

If the bubble subdivides an old edge, the coefficient of
\(\epsilon\) in the product of its two edge factors is
\begin{equation}
 \frac{\eta(B_x+1)(1-A_{x'})
       +\eta'(B_{x'}+1)(1-A_x)}{xx'}.
 \label{eq:detector-bubble-product}
\end{equation}
Substituting the original edge numerator shows that \eqref{eq:detector-bubble-product}
equals
\begin{equation}
 \frac{U(0,x')-U(x,x')}{x}
 +\frac{U(x,0)-U(x,x')}{x'}.
 \label{eq:detector-edge-difference}
\end{equation}
Thus it also gives precisely the two cotangent pullback
corrections.  Bubbles for distinct slots at a vertex are
disjoint, and these are all possibilities for adding one
mark.  Finally, rational-tails graphs have \(h^1=0\) and
trivial automorphism groups, so neither the factor
\(2^{-h^1}\) nor unnormalized gluing introduces another
coefficient.  This proves the identity.
\end{proof}

\subsection{Annihilating the relation span}

\begin{lemma}
\label{lem:detector-annihilation}
For every \(m\geq0\), \(t_m(\Pspan_{g,m})=0\).
\end{lemma}

\begin{proof}
We induct first on arity.  At a fixed arity, the induction already
kills every nontrivial outer gluing and every relation multiplied
by a tail divisor.  With those terms disposed of, degree and
partition length reduce the basic recipes to the all-ones case,
where an alternating-form calculation finishes the argument.

Use the precise span and operation stability in
Proposition~\ref{prop:known-relations}, as established for the
original relation family in \cite[Proposition~1]{Pixton}.
For a basic recipe relation, use the lexicographic order
\((m,r,\ell(\sigma))\in\mathbb Z_{\geq0}^3\):
first the number of labels, then degree, then the number
of partition parts.  At a fixed arity, all relation spaces
at smaller arities, in every degree, have already been
treated.  At zero labels, every admissible basic relation
has positive degree, whereas \(t_0\) kills positive degrees.

Consider a generator obtained by a nontrivial outer gluing.
If that outer graph is not a rational-tails graph, all its
refinements have zero test.  Otherwise, a relation on a
rational factor has zero top integral, by its established
Chow validity.  A relation on the root is killed by the
smaller-arity induction, because that root has fewer labels.
Integration and coproducts then give zero for the entire
generator.

We will also use the following consequence, valid for any
\(b\in\Pspan_{g,m}\):
\begin{equation}
 t_m(D_Tb)=0.
 \label{eq:detector-tail-ideal}
\end{equation}
Indeed, write this product as the gluing pushforward of the
gluing pullback of \(b\).  The latter belongs to the sum of
the relation ideals in the two factors.  Integrate the
rational factor and use the smaller-arity induction on the
root factor, as above.  The root arity is
\(m-|T|+1<m\).  This argument needs no induction on the
degree of \(b\).

It remains to consider a basic recipe relation.  If
\(a_i=3k+h\), with \(k>0\) and \(h=0,1\), its defining
leg factor shows that the relation is \(\psi_i^k\) times
the degree-\((r-k)\) relation with \(a_i=h\).
The admissibility inequality and parity remain unchanged
after subtracting \(3k\) on both sides.  Expand
\(\psi_i=\bar\psi_i+\sum_{T\ni i}D_T\).
The term containing only \(\bar\psi_i\) uses the lower
triple \((m,r-k,\ell(\sigma))\) and is killed by
\eqref{eq:detector-bars}.  Every other term is \(D_Tb'\)
with \(b'\in\Pspan_{g,m}\), by the ideal property, and is
killed by \eqref{eq:detector-tail-ideal}.  That identity
uses smaller arity only, regardless of the degree of
the remaining multiplier.  We may therefore suppose that all
leg indices are zero or one.  We may also suppose \(r\leq m\),
by Lemma~\ref{lem:detector-cutoff}.

If a leg has index zero, extract its degree and partition
coefficient in Lemma~\ref{lem:detector-unit}.  The pullback
term uses \((m-1,r,\ell(\sigma))\) and is killed by
\eqref{eq:detector-forget}.  A remaining term moves a part
\(\ell\) to the zero leg and has order
\((m,r,\ell(\sigma)-1)\).  If \(\ell=1\), this is already
the required strict decrease.  Every other allowed positive
part is \(\ell=3k+h\) with \(k>0\) and \(h=0,1\).
Reducing that leg first gives a bar-multiplied relation
of order \((m,r-k,\ell(\sigma)-1)\), while its tail terms
are again covered by \eqref{eq:detector-tail-ideal}.
Thus every invocation is well founded.  These are legitimate recipe
relations: moving a part preserves its total weight, hence
both admissibility conditions, and removing the zero leg
does not change them.  Moreover
\(r\leq m\leq3g-3+(m-1)\), since \(g\geq2\), so the
pullback term fits the lower-arity dimension bound.
The identity can be used before cutoff, because the test
commutes with pullback and kills the excessive terms.

Only the all-ones case remains.  Its recipe is symmetric
under permutations of the labels, and the same is true of
its test tensor.  Contracting any one slot with \(1\)
gives the test of a forget pushforward, which vanishes by
operation stability and smaller arity.  Contracting with
\(e\) does also: multiply first by
\(\bar\psi_y/(2g-2)\), use the ideal property, and then
forget that slot.

It suffices, therefore, to test all slots on standard
vectors of \(W\).  In the block expansion \eqref{eq:detector-blocks}, a singleton
has zero pairing with an odd input, a positive power of
\(\lambda\) kills odd inputs, and a block of at least four
odd inputs has zero product.  Only undecorated pair and
triple blocks survive.  If there are \(s\) pairs and \(u\)
triples, then
\begin{equation}
 m=2s+3u,\qquad r=s+2u,\qquad
 s+3u=3r-m\geq g+1+|\sigma|>g.
 \label{eq:detector-lagrangian-count}
\end{equation}
The permutation symmetry of the tensor makes its pairing
an ordinary alternating form on the odd input vectors.
Each symplectic pair uses one dual coordinate of \(W_+\);
each cubic uses three such coordinates.  Thus every
alternated term in \eqref{eq:detector-lagrangian-count} uses more than \(g\) coordinates
from the \(g\)-dimensional space \(W_+^\vee\) and is zero.
All standard input tests vanish, so nondegeneracy of the
pairing proves that the tensor itself is zero.
\end{proof}

\subsection{The nonzero value on the candidate}

We finish the proof of Proposition~\ref{prop:detector}.
Every graph in a bridge correction
\(\beta_J^{\mathrm{st}}\) has nontrivial unmarked
stabilization and is not a rational-tails graph.
Its refinements, pullbacks, and products remain outside
the support of the test.  By
\eqref{eq:detector-pair} and \eqref{eq:detector-diagonals},
the remaining part of each factor of \(Y\) tests as
\[
 \Delta_{ab}(1)\Delta_{ac}(1)=\Delta_{abc}(1).
\]
Pair slot \(j\) with the first-Lagrangian vector
\(a_{i_j}\) prescribed by the color word in
Definition~\ref{def:candidate}.  A block with increasing
colors \(i<j<k\) gives \(\theta_{ijk}\).  The product of
the unsymmetrized factors therefore evaluates, up to one
fixed super sign, to
\[
 \Theta=\prod_{1\leq i<j<k\leq g}\theta_{ijk}\neq0,
\]
with the product ordered as the blocks.

Let \(G_{\mathbf i}\) be the stabilizer of the word.
For \(\tau\in G_{\mathbf i}\), transporting its super
permutation across the pairing of odd input slots
contributes \(\operatorname{sgn}(\tau)\); the input word
itself is unchanged.  This cancels the explicit sign of
that summand in \(Y\).  No reordering of the exterior
generators is required in this comparison: all summands
are compared with the same ordered evaluation.
Every color occurs \(\binom{g-1}{2}\) times, and hence
\begin{equation}
 \bigl\langle t_n(Y),
     a_{i_1}\otimes\cdots\otimes a_{i_n}\bigr\rangle
 =
 \pm\left(\binom{g-1}{2}!\right)^g\Theta\neq0.
 \label{eq:detector-candidate-value}
\end{equation}
Together with Lemma~\ref{lem:detector-annihilation}, this
proves the proposition.
